\section{Residual-boundary signs on decorated walks}
\label{sec:residual-boundary}

TPC--68 proves that target signs freeze on a same-residual edge
\citep{WangTPC68}.  We strengthen the bookkeeping from a pure residual walk
to an arbitrary mixed word and identify the exact surviving sign.

\subsection{The edge identity}

Consider an oriented decorated edge
\[
 \xi=(w_-,w_+)
\]
from row \(m\) to row \(n\), with common complete native key.  Abbreviate
\[
 s_-(\xi)=s(w_-),\qquad s_+(\xi)=s(w_+).
\]
The gauged target sign is
\[
 \varepsilon(\xi)
 :=
 \mu(n(w_-))\mu(n(w_+)).
\]

\begin{lemma}[Source and residual factorization of one edge]
\label{lem:edge-sign-factorization}
On the squarefree carrier,
\begin{equation}
 \boxed{
 \varepsilon(\xi)
 =
 \mu(d_m)\mu(d_n)
 \mu(s_-(\xi))\mu(s_+(\xi)).}
 \label{eq:edge-sign-factorization}
\end{equation}
If \(\xi\) is an \(R\)-edge, this reduces to
\begin{equation}
 \varepsilon(\xi)=\mu(d_m)\mu(d_n).
 \label{eq:R-edge-frozen}
\end{equation}
\end{lemma}

\begin{proof}
Since \(n(w)=p(w)r(w)\) is squarefree and \(p(w)\) is prime,
\[
 \mu(n(w))=-\mu(r(w))
 =-\mu(d_{m(w)})\mu(s(w)).
\]
Multiplying the two endpoint identities proves
\eqref{eq:edge-sign-factorization}.  On an \(R\)-edge the two residual
labels agree, so their Möbius factors square to one.
\end{proof}

\subsection{Open paths and closed walks}

\begin{theorem}[Exact residual-boundary sign]
\label{thm:residual-boundary-sign}
Let
\[
 m_0,m_1,\ldots,m_k
\]
be a decorated row path with edge types
\(\omega_\ell\in\{R,E\}\).  Then
\begin{equation}
 \boxed{
 \prod_{\ell=1}^k\varepsilon(\xi_\ell)
 =
 \mu(d_{m_0})\mu(d_{m_k})
 \prod_{\ell:\omega_\ell=E}
 \mu(s_-(\xi_\ell))\mu(s_+(\xi_\ell)).}
 \label{eq:open-path-boundary-sign}
\end{equation}
If \(m_k=m_0\), the source-row boundary disappears:
\begin{equation}
 \boxed{
 \prod_{\ell=1}^k\varepsilon(\xi_\ell)
 =
 \prod_{\ell:\omega_\ell=E}
 \mu(s_-(\xi_\ell))\mu(s_+(\xi_\ell)).}
 \label{eq:closed-walk-erasure-sign}
\end{equation}
\end{theorem}

\begin{proof}
Apply \cref{lem:edge-sign-factorization} to every edge.  Each intermediate
row gauge \(\mu(d_{m_\ell})\) occurs twice and squares to one.  Only the
endpoint row gauges remain on an open path.  Every \(R\)-edge residual
pair also squares to one, leaving exactly
\eqref{eq:open-path-boundary-sign}.  On a closed path the two endpoint
gauges agree and square to one.
\end{proof}

For a closed decorated walk, build a residual multigraph
\(\mathscr G_E\) as follows.  Its vertices are the integer residual labels
appearing on \(E\)-edge endpoints, and every \(E\)-edge contributes the
unordered residual edge
\[
 \{s_-(\xi),s_+(\xi)\}.
\]
The two endpoints are distinct by definition, although parallel edges are
allowed.  Let \(\deg_E(s)\) be the degree of residual vertex \(s\), counted
with multiplicity.

\begin{corollary}[Odd residual boundary]
\label{cor:odd-residual-boundary}
The closed-walk target sign is
\begin{equation}
 \boxed{
 \prod_s\mu(s)^{\deg_E(s)}
 =
 \prod_{\deg_E(s)\ {\rm odd}}\mu(s).}
 \label{eq:odd-residual-boundary}
\end{equation}
In particular, if \(\mathscr G_E\) is Eulerian, the target sign is forced
to be \(+1\).
\end{corollary}

\begin{proof}
Group the factors in
\eqref{eq:closed-walk-erasure-sign} by residual label and use
\(\mu(s)^2=1\).
\end{proof}

\begin{remark}[Edge-count parity is not residual parity]
\label{rem:edge-count-not-enough}
The number of \(E\)-edges does not determine the sign.  Two edges
\((a,b),(b,c)\) leave \(\mu(a)\mu(c)\), while three edges forming the
residual triangle \((a,b),(b,c),(c,a)\) force sign \(+1\).  A nonempty odd
boundary means only that the sign is not algebraically forced by the
interface; its value on one actual carrier can still happen to equal
\(+1\).
\end{remark}

\subsection{The literal closed-walk expansion}

For a closed decorated word of length \(k\), the product of the
normalizing denominators is
\[
 \prod_{\ell=0}^{k-1}W_{m_\ell}.
\]
Combining this with \cref{thm:residual-boundary-sign} gives
\begin{equation}
\boxed{
\begin{aligned}
\tr(H_{\omega_1}\cdots H_{\omega_k})
={}&
\sum_{\substack{m_k=m_0\\
                 \xi_\ell\in
                 \cC^{\omega_\ell}_{m_{\ell-1}m_\ell}}}
\left[
\prod_{\ell:\omega_\ell=E}
\mu(s_-(\xi_\ell))\mu(s_+(\xi_\ell))
\right]
\\
&\times
\frac{
\displaystyle\prod_{\ell=1}^k
\overline{\beta_{w_{\ell,-}}}\beta_{w_{\ell,+}}
}{
\displaystyle\prod_{\ell=0}^{k-1}W_{m_\ell}
}.
\end{aligned}}
\label{eq:literal-residual-boundary-walk}
\end{equation}
Here
\(\cC^R_{mn}\) and \(\cC^E_{mn}\) are the complete equal- and
unequal-residual native collision lists.  The sum retains the native key,
side, block, opposite row, \(\gamma\), time, mask, and every literal
coefficient on every edge.

The complex phase product in
\eqref{eq:literal-residual-boundary-walk} does not generally telescope.
It becomes positive on an exactly atom-consistent backtracking decoration,
but adjacent matrix edges share only a source-row index, not necessarily
an atom.  Treating every complex weight as a vertex gauge would therefore
change the sum.

\begin{corollary}[One-erasure sign]
\label{cor:one-erasure-sign}
Every decorated term in the one-erasure sector has precisely the residual
sign
\begin{equation}
 \mu(s_-(\xi_E))\mu(s_+(\xi_E)),
 \label{eq:one-erasure-sign}
\end{equation}
where \(\xi_E\) is its unique \(E\)-edge.  No same-residual edge contributes
a target Möbius sign after the closed product is formed.
\end{corollary}
